\begin{chapterbox}
    \chapter{Die Hammeraxt (WIP)}
    \label{wip-die-hammeraxt}
    \az{Jahr 58}

    Der Schildzwerg Xoll und der Fährtenleser Vann könnten kaum verschiedener sein. Doch die beiden Väter müssen ihre Kräfte zusammenschließen, um dem Skralvater Hark zu entkommen.
\end{chapterbox}



Diese Fan-Geschichte ist noch in Arbeit. Ich teile sie bereits hier, damit sie für mich nicht verloren geht, falls mein Laptop spontan verschwinden sollte, und damit sie für euch nicht völlig verloren geht, falls ich spontan verschwinden sollte. Fühlt euch in letzterem Fall vollkommen frei, hierrin zu schmökern und nach Lust und Laune weiter daran zu werkeln. Aber solange ich ein aktiver Andori bleibe, lest bitte nicht zu genau nach, was ich geplant haben könnte. LG BBB :D



\newpage
\section{Kapitel}
\az{Jahr 58}

Es war Herbst im Lande Andor. Prinz Thorald stellte einen Spähtrupp gegen die Skrale des Skralbergs zusammen.
\sidenote{\siehechronik{Ein Rabe entscheidet sich (2019)}: Ich war im letzten Herbst von Thorald, des Königs Sohn, als Fährtenleser für einen Spähtrupp gegen die Skrale angeworben, den der Prinz persönlich anführen wollte. Es dauerte Tage, bis der Prinz nüchtern genug war, allein in den Sattel zu kommen. Und ich hatte Zeit, mir die Verteidigungsanlagen in Ruhe anzusehen.}



Vann und Fenn entdeckten eine Höhle im Südlichen Wald und fanden dort Unterschlupf.

Die Schattenhexe Shan kehrte wohlgenährt in ihre Höhle zurück und fand Eindringlinge vor. Xoll vertrieb sie dank eines Edelsteins in seiner Tasche. 

Als die Gruppe die Höhle tiefer erforschten, fanden sie Zeichen, dass die Schattenhexe hier lange genistet hatte. Mithilfe von Xoll und den Edelsteinen machten sie sie fortan sicher vor lichtempfindlichen Schatten.

Xolls Sohn Kjall kam ins Spiel, ohne der Schattenhexe persönlich begegnet zu sein.

Der Trupp wurden in eine Schlacht gegen die Armee des Riesenskrals Hark verwickelt. Menschen und Zwerge kämpften Seite an Seite.
\sidenote{\siehechronik{Ein Rabe entscheidet sich (2019)}: Es gibt viele gute Kämpfer unter den Andori. Wir hatten hier auch eine Schlacht. Und da standen Menschen und Zwerge Seite an Seite.}




\newpage
\section{Epiloge}
\az{Jahr 59}




Vann erschien verletzt in einem Dorf im Rietland, von seinem eigenen Kind besiegt. Fieber plagte ihn. Seine verwundete Schulter war entzündet und eitrig.

Der Händler Nader verhandelte für ihn, der kaum mehr sprechen konnte.

Es gab nur eine alte Heilerin hier im Dorf. Sie sprach von Amputation. Vann vertraute der Heilkunde von Frauen nicht, protestierte, er wolle weiterziehen.

Nader machte ihm klar, dass er auf einem Scheideweg stand. Er konnte weiter über Fenn und Frauen jammern und weiterziehen und elend verenden, oder er konnte die Zähne zusammenbeißen, die dargebotene Hilfe akzeptieren und sich bessern.



\trennzeichen

Fenn träumte von eigenen heroischen Taten.  Schuldgefühle für Vanns Verletzung bissen sich mit der Erleichterung, endlich frei von so vielen Komplikationen in seinem Leben zu sein und seinen eigenen Pfad wählen zu können.

Er reiste weg vom ihm wohlbekannten Südlichen Wald, zum Wachsamen Wald, um zu sehen, ob die \gend{Bewahrer*innen}{Bewahrer} einen Fährtenleser nutzen konnten.


\trennzeichen

Hark starb durch Thorns Hand.

\trennzeichen

Sein Sohn Shron schwor Rache.



\trennzeichen

Xoll Hammeraxt starb bei der Befreiung der Rietburg.

\trennzeichen


Kjall kehrte später zurück, um Xolls Leichnam zu bergen.






